\section{Discussion}
\label{sec:discussion}

\begin{table}[t]
    \centering
    \caption{Verdicts on the five hypotheses fixed before the experiments.}
    \label{tab:verdicts}
    \small
    \begin{tabular}{@{}lp{0.23\textwidth}p{0.17\textwidth}p{0.47\textwidth}@{}}
        \toprule
        & \textbf{Claim} & \textbf{Verdict} & \textbf{Evidence} \\
        \midrule
        H1 & Consistent bias toward A/B & {\bf Not supported} & Position effects are real in 6/7 models but point in different directions; on \mmlu 5/7 are below their average at A. \\
        H2 & 3--8 points inflation on 4-option items & {\bf Not supported} & As published $+$1.3 \ci{0.7}{2.0}, none of it from key placement; A/B key $-$0.7; best case 2.0--3.6. \\
        H3 & Strongest at medium difficulty & {\bf Partly supported} & Strongest where the model is near 50\% accuracy; this is the hard-to-middle range, and the middle third is weaker than the hard third. \\
        H4a & True/false $<$2 points & {\bf Not supported} & Mean spread 14.2 (words) and 4.4 (letters); under 2 for 1/7 and 4/7 models. \\
        H4b & 5-option $>$10 points, the largest & {\bf Partly supported} & Spread $>$10 for 3/7 models, mean 8.8; larger than 4-option on the same questions, but smaller than true/false words and numeric labels. \\
        H5 & Up to 15\% of accuracy near the boundary & {\bf Best case only} & 13--19\% in the hardest strata if every key sat at the model's favored position; 7\% or less as published. \\
        \bottomrule
    \end{tabular}
\end{table}


\Tabref{tab:verdicts} summarizes the verdicts on the five hypotheses.

{\bf Why is ``inflation'' mostly the wrong word?}
A position preference only raises a score when the benchmark's keys are concentrated where the model likes to answer.
\mmlu and \csqa have close-to-uniform keys, so the preferences largely cancel.
What the preferences do produce is \emph{instability}: about 30\% of \mmlu questions and 28\% of \csqa questions are answered correctly under some option orders and incorrectly under others.
This matters for per-question analyses and for small benchmarks, and it would matter for aggregate scores on any benchmark with a skewed key.

\para{Surprises.}
First, the simplest format produced the largest effect: word-answer true/false gives a 14.2-point spread on average.
Second, position A is a disadvantage, not an advantage, for five of these seven models on \mmlu.
Third, the published \mmlu order is 1.3 points easier than its rotations for reasons unrelated to where the key is.
This is a caution for any shuffle-based debiasing: rotating options also damages order-dependent items.

\para{Relation to prior work.}
The model-specific direction, the dependence on label tokens, and the growth of the effect with uncertainty are in line with the findings of \citet{zheng2024selectors} for earlier models.
This comparison is qualitative: we did not rerun prior methods, and we did not conduct a systematic literature review for this study.

\para{Implications for evaluation practice.}
\begin{itemize}[leftmargin=*,itemsep=0pt,topsep=0pt]
    \item Report position-balanced accuracy ($k$ rotations) when comparing models whose scores differ by a few points.
    \item Check the answer-key distribution of any new benchmark; a skewed key is where a 2--5 point distortion would actually appear.
    \item Do not assume two-option or true/false items are immune; test both orders.
    \item Treat per-question correctness as order-dependent for roughly a quarter to a third of items.
\end{itemize}

\subsection{Limitations}
\label{sec:limitations}

\para{Model coverage.}
We study seven open-weight models, none above 14B parameters, five from one vendor.
We planned to run hosted frontier models and could not.
Larger models may be less order-sensitive: within Qwen2.5 the true/false word effect shrinks from 51 points at 1.5B to 4--7 points at 7B and 14B.
Our numbers should not be read as applying to frontier models.

\para{Rotations, not full shuffles.}
Cyclic rotation balances key position with $k$ prompts per question but covers $k$ of $k!$ orders, and it moves order-dependent options out of place.
About 7\% of our \mmlu sample is affected; we did not validate by hand the regular expression that flags these questions.

\para{Position versus label.}
Position and label are tied in every condition except the single A--D versus 1--4 comparison.
We show that the effect is label-dependent but do not apportion it between the two.

\para{First-token read-out.}
Agreement with generated text is 97.5--99.8\% where the model writes a letter, but qwen3-8b and, with numerals, phi3-mini do not reliably start with a label.

\para{One prompt template, zero-shot.}
We did not test few-shot prompts, other instructions, or chain-of-thought answers, all of which may change selection behavior.

\para{True/false is one dataset with one unusual phrasing.}
Our \boolq items are passage-based and use the statement ``The answer to the question \ldots\ is yes.''
The large word-order effect may not carry over to plain factual true/false statements.

\para{Contamination.}
\mmlu, \csqa, and \boolq are very likely in these models' training data.
The 1.1-point residual advantage of the published \mmlu order may reflect this.

\para{Best-position figures are biased upward.}
We choose the best position from the same data.
The noise-only reference (0.6--0.9 points on the full sets) is larger in the smaller difficulty strata.

\para{Difficulty from other models.}
Our difficulty measure is a property of this particular set of seven models and is correlated across the five Qwen models.

\para{Post hoc analyses.}
The decomposition of the as-published gain, the order-dependent split, the split-half replication, the permutation null, and the low-label-mass filter were added after we saw the results.
